\practicechapter{第 5 章实践：编译器 IR、优化 pass 与汇编证据链}

本章对应第一册第 5 章“编译器 IR、优化 pass 与汇编证据链”。实践目标不是要求你实现编译器，而是训练一条判断路径：源码改动或优化等级改变后，二进制证据是否真的变化，变化出现在 section、symbol、反汇编还是 benchmark 数字里。

\practicesection{一、对应原书问题：优化不是一句 Release 更快}

同一份 \code{hello.cpp} 用 \code{-O0} 和 \code{-O2} 编译，得到的二进制可能大小不同，符号数量不同，函数是否保留独立符号也不同。初学者常见错误是直接说“Release 更快”，但没有证明当前跑的就是 Release 产物，也没有证明热函数还以独立形态存在。

实践卷要求先固定产物身份，再讨论优化。固定身份靠 byte size 和 checksum；观察结构靠 section 与 symbol；观察机器码靠 \cmd{objdump -d}；讨论运行时间前，还要知道 benchmark 计时边界里有没有 I/O、打印、构造输入等无关成本。

\practicesection{二、从特例推到规律：函数可能消失}

特例一：\code{int add_i32(int,int)} 在 \code{-O0} 下通常会保留一个可见函数体，反汇编中能找到单独标签。特例二：在 \code{-O2} 下，如果 \code{add_i32} 很小且只被 \code{main} 调用，编译器可能把它 inline 到调用点，符号表中不再以你期待的方式出现。此时不能说“解析器没找到所以函数不存在于源码”，只能说“当前构建产物里没有以独立符号形式保留它”。

由此推出规律：源码层概念和二进制层证据不是一一等价。编译器 IR 与优化 pass 会改变函数边界、常量表达、循环形态和调用结构。实践报告必须写清楚观察层级：源码函数、IR 中的函数、目标文件符号、最终可执行文件符号、反汇编标签，分别可能不同。

\practicesection{三、实践任务：比较两个构建产物}

执行两组构建，并保存输出：

\begin{lstlisting}[language=bash]
g++ -std=c++20 -O0 -g hello.cpp -o hello-o0
g++ -std=c++20 -O2 hello.cpp -o hello-o2

cpu1ee=books/cpu-volume-1-practice/build/executable-evidence-debug/labs/executable_evidence/cpu1ee_cli
$cpu1ee ./hello-o0 > reports/hello-o0.report
$cpu1ee ./hello-o2 > reports/hello-o2.report
$cpu1ee ./hello-o0 --sections-csv > results/hello-o0.sections.csv
$cpu1ee ./hello-o2 --sections-csv > results/hello-o2.sections.csv

objdump -d ./hello-o0 > reports/hello-o0.disasm
objdump -d ./hello-o2 > reports/hello-o2.disasm
\end{lstlisting}

比较时不要只看文件大小。至少检查三类证据：\code{checksum} 是否不同，\code{.text} 的 size 是否变化，反汇编里 \code{add_i32} 是否保留独立标签。若符号没有出现，继续用 \cmd{nm -C} 或 \cmd{objdump -t} 辅助观察，并在报告中写明“未找到独立符号”的范围。

\practicesection{四、验收：优化结论必须降级}

本章交付物是 \filepath{reports/ch05-optimization-evidence.md}。通过标准如下：

\begin{enumerate}
  \item 报告同时包含 \code{-O0} 和 \code{-O2} 的构建命令、checksum 和 byte size。
  \item 至少比较一个 section 字段和一个 symbol 或反汇编证据。
  \item 如果函数被 inline 或符号被优化掉，报告必须写成“当前产物未以独立符号保留”，不能写成“函数不存在”。
  \item 能解释为什么 benchmark 数字必须绑定当前 checksum。
  \item 能指出本卷解析器没有解析 IR，因此 IR 结论需要来自编译器输出或其他工具，而不是从 ELF 字段硬推。
\end{enumerate}

这章完成后，读者应该形成一个习惯：优化相关判断先问“我观察的是哪一层”，再问“这层证据能支持多强的结论”。
